\documentclass[
    language=english,
    doctype=homework,
    institution=none,
]{../../omnilatex}

\RequirePackage{config/document-settings}

\showsolutions

\begin{document}

\maketitle

\begin{exercise}[5]
Define a \emph{singly linked list} and state two advantages it has over a
dynamic array.
\end{exercise}
\begin{solution}
A singly linked list is a linear data structure in which each element (node)
contains a data field and a reference (pointer) to the next node in the
sequence.

Two advantages over a dynamic array:
\begin{enumerate}
    \item $O(1)$ insertion and deletion at the head (no shifting of elements).
    \item No need to pre-allocate or resize a contiguous block of memory.
\end{enumerate}
\end{solution}

\begin{exercise}[10]
Compute the time complexity of the following function using Big-O notation.
\begin{verbatim}
  function mystery(n):
      x = 0
      for i = 1 to n:
          for j = 1 to i:
              x = x + 1
      return x
\end{verbatim}
\end{exercise}
\begin{solution}
The inner loop runs $i$ times for each value of $i$ from $1$ to $n$.
The total number of iterations is:
\[
    \sum_{i=1}^{n} i = \frac{n(n+1)}{2}
\]
Therefore the time complexity is $O(n^{2})$.
\end{solution}

\begin{exercise}[15]
Prove by induction that a binary tree of height $h$ has at most $2^{h+1}-1$
nodes.
\end{exercise}
\begin{solution}
\textbf{Base case} ($h=0$): A tree of height $0$ is a single node.
$2^{0+1}-1 = 1$. \checkmark

\textbf{Inductive step}: Assume a binary tree of height $h$ has at most
$2^{h+1}-1$ nodes. Consider a tree of height $h+1$. Its left and right
subtrees each have height at most $h$, so by the inductive hypothesis each
has at most $2^{h+1}-1$ nodes. Adding the root:
\[
    1 + 2\bigl(2^{h+1}-1\bigr) = 1 + 2^{h+2}-2 = 2^{h+2}-1
\]
This completes the induction. \qed
\end{solution}

\begin{exercise}[10]
Write pseudocode for an in-order traversal of a binary search tree that
collects all values into an array.
\end{exercise}
\begin{solution}
\begin{verbatim}
  function inorder(root, result):
      if root == null:
          return
      inorder(root.left, result)
      append(result, root.value)
      inorder(root.right, result)
\end{verbatim}
For a BST with $n$ nodes this runs in $O(n)$ time and uses $O(h)$ stack
space, where $h$ is the tree height.
\end{solution}

\begin{exercise}[15]
Consider a hash table that uses separate chaining with $m$ slots and $n$
keys. Answer the following:

\begin{enumerate}
    \item[(a)] What is the expected length of a chain under simple uniform
          hashing?
    \item[(b)] If $n = \alpha m$ for load factor $\alpha$, what is the
          expected time for an unsuccessful search?
    \item[(c)] How does open addressing (linear probing) differ in its
          expected search time?
\end{enumerate}
\end{exercise}
\begin{solution}
\begin{enumerate}
    \item[(a)] Under simple uniform hashing the expected chain length is
          $\displaystyle\frac{n}{m} = \alpha$, the load factor.
    \item[(b)] An unsuccessful search must examine the entire chain on
          average, giving $O(1 + \alpha)$ expected time.
    \item[(c)] With linear probing the expected number of probes for an
          unsuccessful search is approximately
          $\dfrac{1}{2}\!\left(1 + \dfrac{1}{(1-\alpha)^{2}}\right)$
          for $\alpha < 1$, which grows faster than separate chaining as
          $\alpha \to 1$.
\end{enumerate}
\end{solution}

\begin{exercise}[5]
State the Master Theorem. You do not need to prove it.
\end{exercise}
\begin{solution}
For a recurrence of the form
$T(n) = a\,T(n/b) + f(n)$
where $a \ge 1$, $b > 1$, and $f$ is asymptotically positive:

\begin{enumerate}
    \item If $f(n) = O(n^{\log_b a - \epsilon})$ for some $\epsilon > 0$,
          then $T(n) = \Theta(n^{\log_b a})$.
    \item If $f(n) = \Theta(n^{\log_b a}\log^{k} n)$, then
          $T(n) = \Theta(n^{\log_b a}\log^{k+1} n)$.
    \item If $f(n) = \Omega(n^{\log_b a + \epsilon})$ and the regularity
          condition holds, then $T(n) = \Theta(f(n))$.
\end{enumerate}
\end{solution}

\end{document}
